\begin{small}
\begin{center}\begin{Huge}R\'esum\'e\end{Huge}\end{center}
\bigskip\bigskip

Nous proposons des m\'ethodes et des outils d'exploration et de
visualisation d'espaces de r\`egles de grande dimension pour des
automates cellulaires. %
Ces objec\-tifs nous ont amen\'e \`a r\'ealiser des exp\'eriences de calcul
parall\`ele cellulaire et \`a concevoir des architectures logicielles bien
adapt\'ees qui nous ont permis de mener \`a bien nos exp\'eriences sur des
machines enti\`erement s\'equentielles.

Nous avons r\'ealis\'e des moteurs cellulaires rapides bas\'es sur une
repr\'esentation des lois par la lookup-table r\'esultante d'un morphisme
et sur une collecte pipeline du voisinage cellulaire. %
Afin de d\'emontrer la validit\'e de cette approche nous avons \'egalement
r\'ealis\'e une impl\'ementation d'un processeur cellulaire SIMD particulier
sous forme du morphisme associ\'e \`a un pr\'ed\'ecodage des instructions pour
une exploitation optimale de la lookup-table; nous proposons en outre
une m\'ethode de transformation d'une lookup-table quelconque en
s\'equence de contr\^ole pour cette machine. %
Nous avons exp\'eriment\'e des modes d'analyse et de visualisation
sp\'ecifiques, en particulier des \aconcept{r\'eductions diff\'erentielles}
et des \aconcept{projections relativistes}. %
Nous avons tourn\'e ces outils vers la production de trajectoires pour
des r\`egles des \aconcept{familles} \arule{Anneal}, \arule{Life},
\arule{Qmean}, \arule{Tube-Worm}, \arule{Fox-Rabbit}, et test\'e les
limites de calcul accessible en r\'ealisant des mesures plus compl\`etes
sur la r\`egle \arule{M46789}. %
En vue de la g\'en\'eralisation de ces exp\'eriences de calcul cellulaire \`a
la visualisation de la structure de l'espace des r\`egles, nous
proposons un nouveau param\`etre de contr\^ole pour l'\'echantillonnage des
espaces de r\`egles, le \aquote{param\`etre d'ench\^assement}, d\'efini pour
chaque \'el\'ement de voisinage du r\'eseau d'interconnexion associ\'e \`a une
r\`egle, comme l'influence de ce voisin sur la r\`egle.

Nous proposons de consid\'erer la programmation cellulaire par tables
comme une m\'ethode d'exploitation des ressources disponibles sur les
machines SISD modernes~(50~Mips, 32~M\'egas), sans pr\'ejudice de l'usage
du calcul cellulaire sur des architectures cellulaires mat\'erielles. %
Nous tentons de d\'efinir les objets \'el\'ementaires permettant de
sp\'ecifier des objets cellulaires arbitrairement complexes constituant
des exp\'eriences de calcul cellulaire. %
Nous consid\'erons d'une part l'existence de dualit\'es \'epist\'emologiques
ou empiriques organisant les objets cellulaires selon des axes
synth\'etique-analytique, temps-espace, algorithme-heuristique,
s\'equentiel-parall\`ele. %
Nous consid\'erons d'autre part la plasticit\'e de l'approche par table,
qui nous permet de projeter les objets du calcul cellulaire
sur ces diff\'erents axes.

Dans le chapitre~\mygetref{space-time-law} nous pr\'esentons la
structure des objets de notre \'etude:~l'espace, le temps et les lois, %
ainsi que les m\'ethodes applicables sur ces objets:~\'edition, r\'eduction,
transformation, visualisation pour l'espace et le temps, ainsi que
construction, alt\'eration, composition, m\'elange et s\'eparation pour les
lois. %
Nous pr\'esentons \'egalement dans ce chapitre l'architecture de nos
moteurs cellulaires. %

Dans le chapitre~\mygetref{rule-space} nous pr\'esentons les outils
conceptuels utilis\'es pour la quantification de l'espace des r\`egles, et
nous proposons un nouveau param\`etre de contr\^ole pour l'exploration et
la structuration de cet espace:~le param\`etre d'ench\^assement. %

Le chapitre~\mygetref{about-anneal} est consacr\'e \`a une \'etude
exp\'erimentale de la r\`egle~\arule{Anneal} bas\'ee sur l'exploitation de nos outils. %
Les exp\'eriences r\'ealis\'ees nous am\`enent \`a proposer l'existence d'une
relative invariance spatio-temporelle des trajectoires associ\'ees \`a
une transition de phase dans l'espace des configurations initiales. %

La conclusion introduit les limites de notre approche qui \`a associ\'e
un constructivisme bottom-up \`a une \aquote{gestion de tables
  g\'en\'eralis\'ee}. %
Nous \'evoquons pour finir une approche alternative de r\'ealisation
cellulaire purement algorithmique. %

\end{small}

%%
%% Local Variables:
%% mode: auto-fill
%% iso-accents-minor-mode: nil
%% TeX-master: "top"
%% TeX-command-default: "mytex"
%% Local IspellParsing: "latex-mode ~latin1"
%% Local IspellPersDict: "~/.ispell_lisp"
%% LocalWords: "CAW Cellular Automata Workbench SIMD SISD Mips M\'egas bottom-up"
%% LocalWords: "morphog\'en\`eses lookup manipulables artefactuelle look-up empty"
%% LocalWords: "impl\'ementation pr\'ed\'ecodage relativistes lookup-table morphisme"
%% LocalWords: "Anneal Qmean Tube-Worm Fox-Rabbit synth\'etique-analytique NEW"
%% LocalWords: "temps-espace algorithme-heuristique s\'equentiel-parall\`ele"
%% LocalWords: "space-time-law rule-space about-anneal constructivisme"
%% End:
